% ============================================================
% APPENDIX
% ============================================================
\appendix

\section{Extended Related Work}
\label{app:related}

\subsection{Adaptive Compute: Reasoning vs.\ Agent Settings}

The adaptive test-time compute literature spans two distinct
settings. In the \emph{reasoning} setting, methods optimize
single-turn chain-of-thought generation on homogeneous benchmarks
(MATH, GSM8K, GPQA).
Signal-based approaches use token-level confidence or entropy to
decide when to extend
thinking~\citep{seag2025, cats2025, corefine2026,
thinkjustenough2025, s1_2025, han2025tokenbudget}.
RL-based methods learn think/no-think policies through reinforcement
learning~\citep{adaptthink2025, thinkless2025, aggarwal2025l1}.
Hidden-state probing methods train lightweight classifiers on
internal representations to estimate task
difficulty~\citep{diffadapt2025, zhu2025llmalreadyknows}.
Because reasoning benchmarks share similar task semantics, signal
direction happens to be consistent across tasks, and the
fixed-direction assumption is never tested.

In the \emph{agent} setting, methods operate in multi-step
interactive environments with external optimizers. Vote-based
methods~\citep{catts2026} sample multiple candidate actions and
trigger evaluation when disagreement is high.
\citet{paglieri2025learning} learn when to invoke a planner at
the episode level. Agentic RL methods~\citep{arpo2025} use
entropy-based adaptive rollout at tool-call steps. All agent-setting
methods inherit the fixed-direction assumption from the reasoning
literature without verifying it. Our work is the first to
systematically study signal semantics across heterogeneous agent
environments, where we find that the direction reverses.

\subsection{Detailed Baseline Comparison}
\label{app:method-comparison}

Table~\ref{tab:baseline-detail} provides a detailed comparison of all
baselines evaluated in this paper. All six methods share a common
assumption: a fixed mapping from signal value to compute need.

\begin{table}[h]
\caption{Detailed comparison of baselines. All methods assume a
fixed signal--utility direction. Phase~1 methods require a separate
calibration dataset; EAAG requires only exploration episodes.}
\label{tab:baseline-detail}
\centering\small
\begin{tabular}{p{1.4cm}p{2.2cm}p{1.8cm}p{2.0cm}p{2.8cm}}
\toprule
Method & Signal & Direction & Granularity & Mechanism \\
\midrule
CATTS & Vote entropy + margin & High disagree. $\to$ trigger
  & Per-step & $K{=}5$ forward passes, arbiter on disagreement \\
SEAG$^\dagger$ & Mean token confidence & Low conf.\ $\to$ search
  & Per-problem & Platt scaling, search depth allocation \\
CaTS$^\dagger$ & Calibrated confidence & Low conf.\ $\to$ generate
  & Per-attempt & Self-calibration, Best-of-N early stopping \\
CoRefine$^\dagger$ & Conv1D confidence & Low conf.\ $\to$ refine
  & Per-problem & Halt / re-examine / redirect controller \\
s1\_budget & Token budget & Fixed budget $\to$ stop
  & Per-problem & Budget-constrained generation \\
AUQ & Verbalized uncertainty & High uncert.\ $\to$ reflect
  & Per-step & Uncertainty-aware memory + reflection~\citep{auq2026} \\
\midrule
\textbf{EAAG} & \textbf{Multi (auto)} & \textbf{Learned}
  & \textbf{Per-step} & Explore + reason + LASSO direction learning \\
\bottomrule
\end{tabular}
\end{table}

\subsection{Adaptive Computation in Neural Networks}

The idea of adapting computation to input difficulty predates LLMs.
Adaptive Computation Time~\citep{graves2016act} allows RNNs to learn
how many computational steps to take per input. Early exit
methods~\citep{teerapittayanon2016branchynet, elhoushi2024layerskip}
attach classifiers at intermediate transformer layers and exit early
for easy inputs. Mixture-of-Experts
architectures~\citep{shazeer2017moe, raposo2024mixtureofdepths}
activate only a subset of parameters per token, adapting capacity
rather than depth. These methods adapt computation within a single
forward pass. Our work operates at a different granularity: we
adapt whether to invoke an external optimizer at each decision step,
and the key challenge is not input difficulty but the
\emph{direction} of the signal--utility relationship.

\subsection{LLM Confidence Estimation and Calibration}

\citet{kadavath2022knowwhat} show that larger models are
well-calibrated on multiple-choice tasks and can estimate P(True)
for their own outputs.
\citet{tao2025revisiting} and \citet{heo2025llmuncertainty} find
that calibration quality varies substantially across task types,
providing converging evidence that uncertainty semantics are
domain-dependent.
\citet{yadkori2024believe} propose decoupling epistemic and
aleatoric uncertainty in LLM outputs, a distinction conceptually
related to our Type~I/Type~D decomposition.
Our work extends these findings: not only does calibration
\emph{quality} vary across tasks, but the \emph{direction} of the
uncertainty--utility relationship reverses entirely.

\subsection{Metareasoning and Value of Computation}

Rational metareasoning~\citep{russell1991right,
horvitz1987reasoning} formalizes the value of computation (VOC)
as a guide for allocating cognitive resources. The classical
framework assumes VOC~$\geq 0$ because an agent can always ignore
computation results. \citet{lieder2020resource} generalize this to
resource-rational analysis, modeling cognition as optimal use of
limited computational resources.
\citet{desabbata2024metareasoning} apply metareasoning to LLMs,
analyzing when to extend chain-of-thought reasoning.
Our work extends this framework by showing that VOC can be
\emph{negative} in the agent setting when the evaluator and executor
are the same model. In this regime, wrong-direction gating does not
merely waste computation (VOC~$= 0$) but actively degrades
performance (VOC~$\ll 0$), violating the classical assumption.
Direction discovery is therefore a necessary precondition for
non-negative VOC across environments.

\subsection{Simpson's Paradox in Machine Learning}

Simpson's paradox~\citep{simpson1951interpretation,
DBLP:books/acm/22/Pearl22n} occurs when a trend present in subgroups
reverses upon aggregation. In machine learning, this phenomenon has
been documented in fairness~\citep{kdd2023simpson}, causal
inference, and treatment effect estimation. Our Two-Source Model
identifies a new instance: within Type~I states, entropy negatively
predicts utility; within Type~D states, entropy positively predicts
utility; but the marginal correlation depends on the mixture
proportion $p_I(\mathcal{E})$, which varies across environments.
This makes direction reversal a structural consequence of
aggregating heterogeneous uncertainty sources.

\subsection{Concurrent Work Statement}

Several 2025--2026 papers independently explore adaptive compute.
CATTS~\citep{catts2026} (Feb 2026) and CoRefine~\citep{corefine2026}
(Feb 2026) are concurrent with no methodological overlap.
AdaptThink~\citep{adaptthink2025} (May 2025) and
Thinkless~\citep{thinkless2025} (May 2025) predate our work but
operate in the reasoning setting. DiffAdapt~\citep{diffadapt2025}
(Oct 2025) assumes a universal U-shaped entropy pattern that our
cross-environment evidence disproves.
AUQ~\citep{auq2026} (Jan 2026) uses verbalized uncertainty for
agent reflection but assumes fixed direction.
Our distinct contribution is threefold: first to (a)~demonstrate
direction reversal across environments and backbones, (b)~formalize
the mechanism via a Two-Source Model grounded in Simpson's paradox,
and (c)~prove that direction discovery is a necessary condition for
cross-environment non-negative VOC.


\section{Experimental Details}
\label{app:experiments}

\subsection{Environment Specifications}

For each of the 8 environments, we provide:
\begin{itemize}[nosep,leftmargin=*]
\item Task description and action space
\item State representation and reward function
\item Episode length distribution
\item Optimizer $T$ implementation details:
  \begin{itemize}[nosep]
  \item HotpotQA, FEVER, TWExpress, Plancraft: per-action evaluation
    (generate $K{=}5$ candidates, evaluate each, select best)
  \item APPS Intro, APPS Interview, CRUXEval: $K$-variant sampling
    (generate $K{=}3$ code solutions, test each, select passing one)
  \item WebShop: LLM-Propose-$K$ (LLM proposes $K{=}3$ action
    sequences, simulate each, select highest-reward)
  \end{itemize}
\end{itemize}

\subsection{Hyperparameter Configuration}

\paragraph{EAAG hyperparameters.}
$N_{\text{explore}} = 50$ episodes,
$\varepsilon_{\text{explore}} = 0.5$,
$\varepsilon_{\text{deploy}} = 0.1 \to 0$ over 100 episodes,
LASSO $\alpha$ selected via 5-fold CV,
gate threshold $\tau = 0.5$,
retrain interval = 30 episodes.

\paragraph{Baseline hyperparameters.}
CaTS: Platt scaling on Phase~1 data (200 episodes), threshold = 0.5.
SEAG: mean confidence threshold calibrated on Phase~1.
CoRefine: entropy threshold calibrated on Phase~1.
CATTS: $K{=}5$ votes, entropy threshold = 0.8, margin threshold = 0.2.
BSW: EAAG with intentionally reversed direction (sign flip).

\subsection{Feature Selection Details}
\label{app:features}

{\color{blue} \textbf{[DATA NEEDED: Table A1.]} Environment $\times$ Feature
matrix showing LASSO coefficients. Rows = 8 environments.
Columns = step\_count, token\_entropy, state\_length, action\_entropy,
num\_available\_actions, has\_output, + LLM-generated features per env.
Values = LASSO coefficient (0 = not selected, sign = direction).
This is the detailed version of Figure 2 in the main text.}

{\color{blue} \textbf{[DATA NEEDED: Table A2.]} LLM-generated features per
environment. Columns: Environment, LLM Prompt (abbreviated), Generated
Features, LASSO Selected (Y/N), Coefficient. Shows the LLM reasoning
step's concrete output for each environment.}

\subsection{Computational Cost Breakdown}

{\color{blue} \textbf{[DATA NEEDED: Table A3.]} Method $\times$ Cost component.
Columns: Method, Phase-1 Episodes, Phase-1 Compute (GPU-hrs),
Per-Step Overhead (ms), Gate Training Time, Total Cost for 500 Episodes.
Shows EAAG's cost advantage: 50 exploration episodes, $\sim$2 GPU-hrs,
0ms gate overhead, $<$1s training.}


\section{Proofs}
\label{app:proofs}

\subsection{Proof of Proposition~\ref{prop:necessity}
  (Necessity of Direction Discovery)}

\begin{proposition*}[Restated]
Suppose there exist two environments $\mathcal{E}_1, \mathcal{E}_2$
with opposite true directions: $d^*(\mathcal{E}_1) = +1$ and
$d^*(\mathcal{E}_2) = -1$. Then no fixed-direction gate $g_d$ can
simultaneously satisfy
$\mathrm{SR}(g_d, \mathcal{E}_1) \geq \mathrm{SR}(\mathrm{base},
\mathcal{E}_1)$ and
$\mathrm{SR}(g_d, \mathcal{E}_2) \geq \mathrm{SR}(\mathrm{base},
\mathcal{E}_2)$.
\end{proposition*}

\begin{proof}
Define a fixed-direction gate as $g_d(s) = \mathbf{1}[d \cdot \sigma(s) > \theta]$
for $d \in \{+1, -1\}$ and threshold $\theta$.

Define $d^*(\mathcal{E}) = +1$ if $\mathrm{Corr}(\sigma(s), U(T,s)) > 0$
in $\mathcal{E}$ (Type~D), and $d^*(\mathcal{E}) = -1$ if
$\mathrm{Corr}(\sigma(s), U(T,s)) < 0$ (Type~I).

\textbf{Lemma (Wrong-direction harm).} If $d \neq d^*(\mathcal{E})$, then
$\mathrm{SR}(g_d, \mathcal{E}) \leq \mathrm{SR}(\mathrm{base}, \mathcal{E}) - \delta(\mathcal{E})$
where $\delta(\mathcal{E}) > 0$ depends on signal strength $|\rho(\mathcal{E})|$.

\emph{Proof of Lemma}: By definition of $d^*$, states with high
$d \cdot \sigma(s)$ are states where $U < 0$. The gate triggers $T$
at these states, causing the agent to adopt worse actions. The base
policy, which never triggers $T$, avoids this systematic harm.

\textbf{Main argument:}

\emph{Case 1}: $d = +1$. In $\mathcal{E}_2$ where $d^* = -1$,
$d \neq d^*(\mathcal{E}_2)$. By Lemma:
$\mathrm{SR}(g_{+1}, \mathcal{E}_2) \leq \mathrm{SR}(\mathrm{base}, \mathcal{E}_2) - \delta(\mathcal{E}_2)
< \mathrm{SR}(\mathrm{base}, \mathcal{E}_2)$.
Constraint violated.

\emph{Case 2}: $d = -1$. In $\mathcal{E}_1$ where $d^* = +1$,
by symmetric argument, constraint violated.

Since $d \in \{+1, -1\}$ and both cases lead to violation, no fixed
$d$ satisfies both constraints simultaneously.

\textbf{Empirical grounding:}
$\delta(\text{HotpotQA}) \approx 37.0$\,pp (BSW ablation),
$\delta(\text{WebShop}) \approx 23.2$\,pp,
$\delta(\text{FEVER}) \approx 36.8$\,pp.
MLP with wrong direction: 45.3\% $<$ base 49.0\% on HotpotQA.
\end{proof}

\subsection{Wrong-Direction Damage Quantification}

\begin{table}[h]
\centering\small
\caption{Wrong-direction damage scales with signal strength.}
\begin{tabular}{lcccc}
\toprule
Environment & Correct SR & Wrong SR & $\Delta$ SR & $|\rho|$ (strongest) \\
\midrule
HotpotQA   & 95.2\% & 58.2\% & $-$37.0 & 0.494 \\
WebShop    & 43.8\% & 21.4\% & $-$22.4 & 0.444 \\
FEVER      & 49.8\% & 13.0\% & $-$36.8 & 0.619 \\
APPS Intro & 66.0\% & 62.5\% & $-$3.5  & 0.155 \\
\bottomrule
\end{tabular}
\end{table}

The damage magnitude scales with signal strength: environments with
stronger signal--utility coupling ($|\rho|$) suffer larger SR degradation
when direction is reversed ($R^2 > 0.5$). On HotpotQA, the MLP gate
with wrong direction achieves 45.3\%, \emph{below} the 49.0\% base,
demonstrating that a more powerful gate with wrong direction is
\emph{worse} than a simpler gate with wrong direction (logistic: 62.0\%),
because it more precisely targets harmful states.

\subsection{VOC Non-Negativity Scope}

In the classical VOC framework~\citep{russell1991right},
$\mathrm{VOC}(T, s) = \mathbb{E}[\max(V(a_T), V(a_{\mathrm{base}}))] - V(a_{\mathrm{base}}) \geq 0$
because the agent can \emph{disregard} the computation result if
unfavorable. Our setting violates this: when the agent uses an
evaluator/optimizer $T$, the evaluator's assessment \emph{is} the
agent's decision. The agent cannot separately evaluate the
evaluator's output because the evaluator \emph{is} the evaluation
mechanism.

Under evaluator-executor identity, VOC can be negative:
$\mathrm{VOC}(T, s) = \mathbb{E}[V(a_T(s))] - V(a_{\mathrm{base}}(s))$.
If $T$ systematically selects worse actions in Type~I states,
$\mathrm{VOC} < 0$ for those states. This explains why wrong-direction
gating is not merely wasteful ($\mathrm{VOC} = 0$) but actively
harmful ($\mathrm{VOC} < 0$).


\section{Two-Source Model: Full Derivation and Verification}
\label{app:model}

\subsection{Full Derivation}

At each step $t$, state $s_t$ belongs to one of two latent types:
$s_t \in \text{Type~I}$ with probability $p_I(\mathcal{E})$,
$s_t \in \text{Type~D}$ with probability $1 - p_I(\mathcal{E})$.

Conditional utility models:
\begin{align}
  U(T, s) \mid s \in \text{Type~I} &= -\alpha \cdot H(s) + \varepsilon_I, \quad \alpha > 0 \\
  U(T, s) \mid s \in \text{Type~D} &= +\beta \cdot H(s) + \varepsilon_D, \quad \beta > 0
\end{align}

The marginal correlation:
\begin{align}
  \mathrm{sign}(\rho(\mathcal{E})) &= \mathrm{sign}(\beta \cdot (1 - p_I) - \alpha \cdot p_I) \\
  &= \mathrm{sign}(\beta - (\alpha + \beta) \cdot p_I)
\end{align}

Reversal threshold: $p_I^* = \beta / (\alpha + \beta)$.

\subsection{Environment Mapping on the $p_I$ Axis}

{\color{red} \textbf{[FIGURE NEEDED: Figure A1.]} 1D axis diagram.
Horizontal axis from $p_I = 0$ (pure Type~D) to $p_I = 1$ (pure Type~I).
Mark each environment's estimated position with a labeled point.
Show the reversal threshold $p_I^*$ as a vertical dashed line.
Left of $p_I^*$: positive $\rho$ (green). Right: negative $\rho$ (red).
Provides a visual summary of the entire Two-Source Model.}

Estimated positions on the $p_I$ axis (0 = pure Type~D, 1 = pure Type~I):
\begin{itemize}[nosep,leftmargin=*]
\item $p_I \approx 0.2$: APPS Interview (strong positive $\rho = +0.317$)
\item $p_I \approx 0.45$: APPS Intro (near-zero $\rho = +0.012$)
\item $p_I \approx 0.5$: WebShop (entropy $\rho \approx 0$, but non-entropy signals work)
\item $p_I \approx 0.55$: CRUXEval (weak negative $\rho = -0.048$)
\item $p_I \approx 0.6$: Plancraft (weak negative, rollouts harmful)
\item $p_I \approx 0.7$: HotpotQA (negative $\rho = -0.041$)
\item $p_I \approx 0.8$: TWExpress (negative $\rho = -0.290$)
\item $p_I \approx 0.9$: FEVER (strong negative $\rho = -0.119$)
\end{itemize}

\subsection{Prediction Verification: Full Data}

{\color{blue} \textbf{[DATA NEEDED: Table A4.]} P1 Temporal Dynamics full data.
Columns: Environment, $\rho_{\text{early}}$ (step $\leq$ median),
$\rho_{\text{late}}$ (step $>$ median), $p_{\text{early}}$, $p_{\text{late}}$.
All 8 environments. Shows the temporal shift pattern for each.}

{\color{blue} \textbf{[DATA NEEDED: Table A5.]} P3 Signal Identity Alignment.
Columns: Environment, Dominant Type, Strongest Signal, $|\rho|$,
Signal Interpretation. Shows what each environment's strongest
signal measures and how it aligns with the predicted type.}


\section{Multi-Backbone Verification}
\label{app:multi-backbone}

We validate our core findings on three backbones from three vendors:
Qwen3-4B-Instruct (Alibaba, 4B), Phi-3.5-mini-instruct (Microsoft, 3.8B),
and Llama-3.1-8B-Instruct (Meta, 8B) across all 8 environments.

\subsection{Signal Direction: Environment $\times$ Model Interaction}

The multi-backbone results (Table~\ref{tab:signal-discovery} in main text)
reveal that 4/5 environments with $\geq$2 significant signals show
cross-model direction variation. Structural signals
(\texttt{step\_count}) are model-invariant while entropy signals are
model-dependent.

{\color{blue} \textbf{[DATA NEEDED: Table A6.]} Full multi-backbone results.
8 environments $\times$ 3 backbones $\times$ multiple signals.
Extends Table 1 in main text with all signal types, not just entropy.
Shows that step\_count direction is stable while entropy flips.}

\subsection{Fixed-Direction Baselines Fail Across Backbones}

On Phi-3.5 APPS: all 6 CB methods score \emph{below} base\_only
(59\% $\to$ 27--36\%), demonstrating that fixed-direction assumption
is catastrophically wrong when the model changes.

{\color{blue} \textbf{[DATA NEEDED: Table A7.]} Cross-backbone baseline results.
Methods $\times$ environments for Phi-3.5 and Llama-3.1 backbones.
Shows that the same baselines that work on Qwen3 fail on other backbones,
while EAAG adapts successfully.}


\section{Additional Analyses}
\label{app:additional}

\subsection{Hyperparameter Sensitivity}

{\color{red} \textbf{[FIGURE NEEDED: Figure A2.]} Hyperparameter sensitivity plots.
(a) SR vs $N_{\text{explore}} \in \{10, 20, 30, 50, 100\}$ for 4 environments.
Shows stability for $N_{\text{explore}} \geq 30$.
(b) SR vs gate threshold $\tau \in \{0.3, 0.4, 0.5, 0.6, 0.7\}$.
Shows robustness across threshold range.
Similar to standard ablation sensitivity plots in NeurIPS papers.}

EAAG is robust to $N_{\text{explore}} \in [30, 100]$ (stable SR for
$N_{\text{explore}} \geq 30$; marginal gains above 50) and gate
threshold $\tau \in [0.3, 0.7]$.

\subsection{Statistical Significance}

All head-to-head comparisons use 3-seed evaluation (200 episodes/seed).
Bootstrap confidence intervals (5000 resamples) are reported for all
comparisons. 18 of 34 wins are statistically significant (95\% CI
excludes zero). Pareto-dominance claims use a conservative criterion
that does not rely on p-values.

{\color{blue} \textbf{[DATA NEEDED: Table A8.]} Full statistical significance table.
For each method $\times$ environment pair: $\Delta$SR, 95\% bootstrap CI,
significant (Y/N). 48 cells total (6 baselines $\times$ 8 environments).}

\subsection{Controlled Information Manipulation: Full Results}
\label{app:controlled}

To provide stronger evidence that information structure drives signal
semantics, we manipulate information availability within HotpotQA
itself. \emph{InfoPoor} (search returns only the first sentence)
reduces base SR to 44.0\% while \emph{InfoRich} (gold evidence
injected at start) raises it to 82.0\%. The signal hierarchy shifts:
in InfoPoor, \texttt{step\_count} dominates ($\rho{=}{-}0.608$) while
entropy carries a weak positive signal ($\rho{=}{+}0.119$,
$n{=}317$; note that the Two-Source Model predicts entropy should be
negative in info-poor settings, but the small sample and mixed
early-step Type~D component may explain the weakly positive sign;
the key finding is that \texttt{step\_count} overtakes entropy as the
dominant signal); in InfoRich, entropy becomes the dominant positive
signal ($\rho{=}{+}0.311$) while \texttt{step\_count} weakens
($\rho{=}{-}0.147$). This demonstrates that the \emph{identity} of
the most informative signal, not just its direction, is a function
of the environment's information structure, supporting Observation~2.

{\color{blue} \textbf{[DATA NEEDED: Table A9.]} InfoPoor vs InfoRich
full signal correlation table for all features. Columns: Condition,
Base SR, $\rho$ for each feature, strongest signal.}

\subsection{LLM Reasoning Prompts and Outputs}
\label{app:llm-prompts}

{\color{blue} \textbf{[DATA NEEDED.]} Full LLM prompt template used in
the Reason step of EAAG. Include one complete example showing:
(1) the prompt given to the LLM with exploration data summary,
(2) the LLM's full output including environment profile and
feature hypotheses, (3) which features LASSO subsequently selected.
Use WebShop as the example (most interesting LLM contribution).
This is important for reproducibility (NeurIPS checklist item).}


\section{Broader Impact}
\label{app:impact}

This work addresses computational efficiency in LLM agent deployment.
By reducing unnecessary optimizer invocations, EAAG lowers the energy
and cost footprint of test-time compute without sacrificing task
performance. The prescriptive framework (characterize information
structure before designing gates) may help practitioners avoid
deploying fixed-direction methods that actively harm performance in
mismatched environments. We do not foresee direct negative societal
impacts beyond those inherent to the LLM agents themselves.
